\qquad A seguir apresenta-se a função responsável por orquestrar, para
cada instância, a geração do vetor, a execução do algoritmo de
ordenação selecionado (recebido como um ponteiro de função, o que
permite reutilizar a mesma rotina para os quatro algoritmos deste
trabalho) com a respectiva medição de tempo, e o registro dos
arquivos de entrada, saída e tempo dentro da pasta correspondente ao
algoritmo — pasta essa que também é criada automaticamente pelo
próprio programa, caso ainda não exista (arquivos
\texttt{base.h} e \texttt{main.cpp}):

\begin{verbatim}
void criarEstruturaPastas(const string &algoritmo) {
    const string grupos[3] = {"Arquivos de Entrada", "Arquivos de Saida",
                               "Arquivos de Tempo"};
    const string categorias[3] = {"Crescente", "Decrescente", "Random"};
    for (const string &grupo : grupos)
        for (const string &categoria : categorias)
            fs::create_directories(algoritmo + "/" + grupo + "/" + categoria);
}

void operacoes(const string &algoritmo, FuncaoOrdenacao funcao,
               int tipo, int tamanho) {
    criarEstruturaPastas(algoritmo);

    clock_t start_t, end_t;
    double tempoGasto;

    int *vetor = gerarSequencia(tipo, tamanho);
    salvarEntrada(algoritmo, tipo, tamanho, vetor);

    start_t = clock();
    funcao(vetor, tamanho);
    end_t = clock();

    tempoGasto = (end_t - start_t) / (double)CLOCKS_PER_SEC;

    salvarTempo(algoritmo, tipo, tamanho, tempoGasto);
    salvarSaida(algoritmo, tipo, tamanho, vetor);

    delete[] vetor;
}
\end{verbatim}
